\documentclass[10pt, letterpaper]{article}

% Inhaltsverzeichnis für Pakettypen (nur für Übersicht im Header, wird nicht im Dokument angezeigt)
% 1. Seitenlayout und Ränder
% 2. Sprache und Zeichensatz
% 3. Mathematik und Theorem-Umgebungen
% 4. Eigene Makros
% 5. Diagramme und Grafiken
% 6. Tabellen und Aufzählungen
% 7. Inhaltsverzeichnis
% 8. Abschnittsüberschriften
% 9. Abstrakt-Umgebung
% 10. Todos/Notizen
% 11. Rahmen/Box-Umgebungen
% 12. Python-Integration
% 13. Literaturverwaltung
% 14. Hyperlinks
% 15. Absatzeinstellungen
% 16. Umgebungen
% 17  Graphik
% 18  Extra
% 00. Titel und Autor

\usepackage{slashed}

% --- 1. Seitenlayout und Ränder ---
\usepackage[margin=3cm]{geometry}

% --- 2. Sprache und Zeichensatz ---
\usepackage[english]{babel}
\usepackage[T1]{fontenc}
\usepackage[utf8]{inputenc}

% --- 3. Mathematik und Theorem-Umgebungen ---
\usepackage{amsmath, amssymb, amsthm}
\usepackage{mathrsfs}
\DeclareMathOperator{\WF}{WF}

% --- 4. Eigene Makros ---
\usepackage{xcolor}
\newcommand{\SKP}{\langle\cdot,\cdot\rangle}
\newcommand{\id}{\operatorname{id}}
\newcommand{\sgn}{\operatorname{sgn}}
\newcommand{\Ad}{\operatorname{Ad}}
\newcommand{\ad}{\operatorname{ad}}
\newcommand{\K}{\mathbb{K}}
\newcommand{\R}{\mathbb{R}}
\newcommand{\N}{\mathbb{N}}
\newcommand{\Q}{\mathbb{Q}}
\newcommand{\Z}{\mathbb{Z}}
\newcommand{\C}{\mathbb{C}}
\newcommand{\QU}{\mathbb{H}}
\newcommand{\Cinf}{C^{\infty}}
\newcommand{\Mod}{\mathrm{Mod}}
\newcommand{\Alg}{\mathrm{Alg}}
\newcommand{\Diff}{\mathrm{Diff}}
\newcommand{\Iso}{\mathrm{Iso}}
\newcommand{\Aut}{\mathrm{Aut}}
\newcommand{\End}{\mathrm{End}}
\newcommand{\Hom}{\mathrm{Hom}}
\newcommand{\Cl}{\mathrm{Cl}}
\newcommand{\entwurf}[1]{\textcolor{red}{#1}}
\newcommand{\GL}{\mathrm{GL}}
\newcommand{\Mat}{\mathrm{Mat}}
\newcommand{\SL}{\mathrm{SL}}
\newcommand{\SO}{\mathrm{SO}}
\newcommand{\SU}{\mathrm{SU}}
\newcommand{\Sp}{\mathrm{Sp}}
\newcommand{\Spin}{\mathrm{Spin}}
\newcommand{\Pin}{\mathrm{Pin}}
\newcommand{\Lor}{\mathrm{Lor}}
\newcommand{\Ogroup}{\mathrm{O}}
\newcommand{\U}{\mathrm{U}}
\newcommand{\CP}{\mathbb{C} P}
\newcommand{\RP}{\mathbb{R} P}

% --- 5. Diagramme und Grafiken ---
\usepackage{graphicx}
\graphicspath{{../../Library/Images/}}
\usepackage[export]{adjustbox}
\usepackage{tikz}
\usetikzlibrary{decorations.pathreplacing, arrows.meta, positioning}
\usepackage{tikz-cd}

% --- 6. Tabellen und Aufzählungen ---
\usepackage{enumitem}
\setlist[itemize]{left=0.5cm}

\newenvironment{romanenum}[1][]
  {%
    \ifx&#1&
    \else
      \textbf{#1}\quad
    \fi
    \begin{enumerate}[label=\roman*)]
  }
  {%
    \end{enumerate}%
  }

% --- 7. Inhaltsverzeichnis ---
\usepackage{tocloft}
\renewcommand{\cftsecfont}{\footnotesize}
\renewcommand{\cftsubsecfont}{\footnotesize}
\renewcommand{\cftsubsubsecfont}{\footnotesize}
\renewcommand{\cftsecpagefont}{\footnotesize}
\renewcommand{\cftsubsecpagefont}{\footnotesize}
\renewcommand{\cftsubsubsecpagefont}{\footnotesize}
\usepackage{etoc}

% --- 8. Abschnittsüberschriften ---
\usepackage{titlesec}
\titleformat{\section}{\normalfont\large\bfseries}{\thesection}{1em}{}
\titleformat{\subsection}{\normalfont\normalsize\bfseries}{\thesubsection}{0.5em}{}
\titleformat{\subsubsection}{\normalfont\normalsize\bfseries}{\thesubsubsection}{0.5em}{}
\setcounter{secnumdepth}{4}

% --- 9. Abstrakt-Umgebung ---
\usepackage{changepage}
\renewenvironment{abstract}
  {
    \begin{adjustwidth}{1.5cm}{1.5cm}
    \small
    \textsc{Abstract. –}%
  }
  {
    \end{adjustwidth}
  }

% --- 10. Todos/Notizen ---
\usepackage{todonotes}
% Hellblau definieren
\definecolor{hellblau}{rgb}{0.3,0.6,1}
\newenvironment{note}{\par\begingroup\color{hellblau}\itshape}{\par\endgroup}

% --- 11. Rahmen/Box-Umgebungen ---
\usepackage{mdframed}
\usepackage{tcolorbox}
\colorlet{shadecolor}{gray!25}

\newenvironment{customTheorem}
  {\vspace{10pt}%
   \begin{mdframed}[
     backgroundcolor=gray!20,
     linewidth=0pt,
     innertopmargin=10pt,
     innerbottommargin=10pt,
     skipabove=\dimexpr\topsep+\ht\strutbox\relax,
     skipbelow=\topsep,
   ]}
  {\end{mdframed}
   \vspace{10pt}%
  }

% --- 12. Python-Integration ---
% (Deaktiviert in dieser Version, aktiviere bei Bedarf)
% \usepackage{pythontex}
% \usepackage[makestderr]{pythontex}

% --- 13. Literaturverwaltung ---
\usepackage{csquotes}
\usepackage[backend=biber, style=alphabetic, citestyle=alphabetic]{biblatex}
\addbibresource{bibliography.bib}

% --- 14. Hyperlinks ---
\usepackage{hyperref}
\hypersetup{
  colorlinks   = true,
  urlcolor     = blue,
  linkcolor    = blue,
  citecolor    = blue,
  frenchlinks  = true
}

% --- 15. Absatzeinstellungen ---
\usepackage[parfill]{parskip}
\sloppy

% --- 16. Umgebungen ---
\usepackage{thmtools}

\newcommand{\CustomHeading}[3]{%
  \par\medskip\noindent%
  \textbf{#1 #2} \textnormal{(#3)}.\enskip%
}

\newenvironment{DEF}[2]{\CustomHeading{Definition}{#1}{#2}}{}
\newenvironment{PROP}[2]{\CustomHeading{Proposition}{#1}{#2}}{}
\newenvironment{THEO}[2]{\CustomHeading{Theorem}{#1}{#2}}{}
\newenvironment{LEM}[2]{\CustomHeading{Lemma}{#1}{#2}}{}
\newenvironment{KORO}[2]{\CustomHeading{Corollar}{#1}{#2}}{}
\newenvironment{REM}[2]{\CustomHeading{Remark}{#1}{#2}}{}
\newenvironment{EXA}[2]{\CustomHeading{Example}{#1}{#2}}{}
\newenvironment{STUD}[2]{\CustomHeading{Study}{#1}{#2}}{}
\newenvironment{CONC}[2]{\CustomHeading{Concept}{#1}{#2}}{}
\newenvironment{OTH}[2]{\CustomHeading{Other}{#1}{#2}}{}
\newenvironment{EXE}[2]{\CustomHeading{Exercise}{#1}{#2}}{}
\newenvironment{MOT}[2]{\CustomHeading{Motivation}{#1}{#2}}{}
\newenvironment{PROOF}[2]{\CustomHeading{Proof}{#1}{#2}}{}

% --- Unit Umgebung für Source-Inhalte ---
\usepackage{mdframed}
\newmdenv[
  linewidth=1pt,
  topline=false,
  bottomline=false,
  rightline=false,
  leftmargin=0cm,
  rightmargin=0cm,
  skipabove=10pt,
  skipbelow=10pt,
  innertopmargin=0.5\baselineskip,
  innerbottommargin=0.5\baselineskip,
  backgroundcolor=gray!10,
  linecolor=gray
]{unitbox}

\newenvironment{unit}[1]
  {\begin{unitbox}\textbf{Unit #1}\par\smallskip}
  {\end{unitbox}}

% --- 17. ... ---

% --- 18. Extras ---
\usepackage{stmaryrd}
\usepackage{bbold}  % falls du \mathbb{1} nutzen willst

\begin{document}

\section{Einführung Riemannsche Spin-Geometrie (Handout)}

\subsection{Bündeltheorie}

\begin{DEF}{}{Faserbündel}
Seien $E$, $M$ und $F$ glatte Mannigfaltigkeiten. Ein \emph{Faserbündel} mit Totalraum $E$, Basisraum $M$ und typischer Faser $F$
ist eine glatte surjektive Abbildung
\[
\pi:E\to M,
\]
so dass es zu jedem $x\in M$ eine offene Umgebung $U\subset M$ und einen Diffeomorphismus
\[
\Phi_U:\pi^{-1}(U)\to U\times F
\]
gibt mit kommutierenden Diagramm
\[
\begin{tikzcd}
\pi^{-1}(U) \arrow[r,"\Phi_U"] \arrow[d,"\pi"'] & U\times F \arrow[dl,"\operatorname{pr}_1"] \\
U &
\end{tikzcd}
\]
Insbesondere ist für jedes $y\in U$ die \emph{Faser} $E_y:=\pi^{-1}(y)$ \emph{über} $y$ isomorph zu $\{y\}\times F\cong F$.
\end{DEF}

\begin{DEF}{}{Vektorbündel}
Ein \emph{Vektorbündel} ist ein Faserbündel $\pi:E\to M$ mit typischer Faser $V$ einem endlichdimensionalen Vektorraum $V$, sodass gilt:
\begin{enumerate}
  \item Jede Faser $E_x=\pi^{-1}(x)$ trägt die Struktur eines Vektorraums.
  \item Die lokalen Trivialisierungen
        \[
        \Phi_U:\pi^{-1}(U)\to U\times V
        \]
        sind faserweise linear.
\end{enumerate}
\end{DEF}

\begin{DEF}{}{Prinzipalbündel}
Sei $G$ eine Lie-Gruppe. Ein \emph{Prinzipalbündel}
mit Strukturgruppe $G$ ist ein Faserbündel $\pi:P\to M$ zusammen mit einer rechten Wirkung
\[
R: P\times G\to P,\qquad (p,g)\mapsto p\cdot g,
\]
so dass gilt:
\begin{enumerate}
    \item die Rechtswirkung \emph{erhält die Faser unter} $\pi$, d.h. 
    \[
    \begin{tikzcd}
    P \arrow[r,"R_g"] \arrow[dr,"\pi"'] & P \arrow[d,"\pi"] \\
    & M
    \end{tikzcd}
    \]
    kommutiert für alle $g\in G$,
    \item die Wirkung ist \emph{frei}, d.\,h.
    \[
    p\cdot g=p \;\Rightarrow\; g=e,
    \]
    \item die Wirkung ist auf jeder Faser $P_x:=\pi^{-1}(x)$ \emph{transitiv}, d.h. für alle $p,q\in P_x$ existiert ein $g\in G$, sodass 
    $$p\cdot g=q$$
    \item lokale Tivialisierungen $\Phi_u:\pi^{-1}(U)\to U\times G$ sind $G$-\emph{äquivariant}, d.h. das Diagramm
    \[
    \begin{tikzcd}
    \pi^{-1}(U)\times G \arrow[r,"\ \ R_G\ "] \arrow[d,"\Phi_U\times \id_G"'] & \pi^{-1}(U) \arrow[d,"\Phi_U"] \\
    (U\times G)\times G \arrow[r,"\ \ r_G\ "] & U\times G
    \end{tikzcd}
    \]
    kommutiert bzw. es gilt $(x,h)\cdot g=(x,hg)$.
\end{enumerate}
\end{DEF}


\begin{DEF}{}{Assoziiertes Vektorbündel}
Sei \(\pi:P\to M\) ein Prinzipalbündel mit Strukturgruppe \(G\) und $\rho:G\to \GL(V)$ eine lineare Darstellung auf einem Vektorraum \(V\). Das \emph{assoziierte Vektorbündel}
ist der Quotient
\[
P\times_\rho V:=(P\times V)/\sim,
\]
wobei
\[
(p,v)\sim (p\cdot g,\rho(g^{-1})v).
\]
Die Projektion
\[
[p,v]\mapsto \pi(p)
\]
macht \(P\times_\rho V\) zu einem Vektorbündel über \(M\) mit typischer Faser \(V\) und Vektorraumstruktur auf jeder Faser durch
\[
[p,v]+[p,w]:=[p,v+w],
\qquad
\lambda [p,v]:=[p,\lambda v].
\]
\end{DEF}



\subsection{Clifford-Algebra von Vektorräumen mit symmetrischen Bilinearformen}

\begin{DEF}{}{Unitale $\mathbb{K}$-Algebra}
Eine \emph{unitale $\mathbb{K}$-Algebra} ist ein $\mathbb{K}$-Vektorraum $A$ zusammen mit
einer bilinearen Multiplikation
\[
A\times A\to A,\qquad (a,b)\mapsto ab,
\]
so dass $A$ mit dieser Multiplikation ein assoziativer Ring ist mit einem Einselement
\[
1_A\in A \quad \text{ mit}\quad 1_Aa=a1_A=a
\qquad\text{für alle }a\in A.
\]
\end{DEF}

\begin{DEF}{}{Algebrenhomomorphismus}
Seien $A$ und $B$ unitale $\mathbb{K}$-Algebren. Eine Abbildung $\varphi:A\to B$
heißt \emph{$\mathbb{K}$-Algebrenhomomorphismus}, falls gilt:
\begin{enumerate}
    \item \(\varphi\) ist \(\mathbb{K}\)-linear,
    \item \(\varphi\) ist multiplikativ, d.h.
    \[
    \varphi(ab)=\varphi(a)\varphi(b)
    \qquad\text{für alle }a,b\in A,
    \]
    \item \(\varphi\) erhält das Einselement, d.h.
    \[
    \varphi(1_A)=1_B.
    \]
\end{enumerate}
\end{DEF}

\begin{DEF}{}{Clifford-Algebra von \((V,\beta)\)}
Sei \(\mathbb{K}=\mathbb{R}\) oder \(\mathbb{C}\), und sei \(V\) ein endlichdimensionaler \(\mathbb{K}\)-Vektorraum, ausgestattet mit einer symmetrischen Bilinearform
\[
\beta:V\times V\to \mathbb{K}.
\]

\begin{enumerate}
\item Eine \emph{Clifford-Algebra} zu \((V,\beta)\) ist eine unitale \(\mathbb{K}\)-Algebra
\(\Cl(V,\beta)\) zusammen mit einer linearen Abbildung
\[
\iota:V\to \Cl(V,\beta),
\]
so dass die folgenden Eigenschaften erfüllt sind:

\begin{enumerate}
    \item Für alle \(v\in V\) gilt
    \[
    \iota(v)^2=-\beta(v,v)\cdot 1.
    \]
    Äquivalent dazu gilt für alle \(v,w\in V\):
    \[
    \iota(v)\iota(w)+\iota(w)\iota(v)=-2\beta(v,w)\cdot 1.
    \]

    \item Das Paar \((\Cl(V,\beta),\iota)\) ist universal bezüglich der Clifford-Relation, d.h.:
    Ist \(A\) eine unitale \(\mathbb{K}\)-Algebra und
    \[
    \iota':V\to A
    \]
    eine lineare Abbildung mit
    \[
    \iota'(v)^2=-\beta(v,v)\cdot 1
    \qquad\text{für alle }v\in V,
    \]
    so existiert genau ein Algebra-Homomorphismus
    \[
    \varphi:\Cl(V,\beta)\to A,
    \]
    so dass das Diagramm
    \[
    \begin{tikzcd}
    V \arrow[r, "\iota"] \arrow[dr, "\iota'"'] & \Cl(V,\beta) \arrow[d, "\varphi"] \\
    & A
    \end{tikzcd}
    \]
    kommutiert.
\end{enumerate}

\item Falls der Vektorraum $V$ reell ist, ist die komplexifizierte Clifford-Algebra von $\Cl_n(V,\beta)$ definiert durch 
  \[
  \C l(V,g):=\Cl(V,g)\otimes_{\R}\C.
  \]
\end{enumerate}
\end{DEF}



\begin{REM}{}{Eigenschaften der Clifford-Algebra}
\begin{enumerate}
  \item \emph{Basis von $\Cl(V,\beta)$}: Sei $V$ ein $\K$-Vektorraum mit Basis $b_1, \ldots, b_n \in V$. Dann ist 
    $$
    \left\{b_{i_1} \cdot \ldots \cdot b_{i_k}\right\}_{\substack{k=0,1, \ldots n \\ 1 \leq i_1<i_2<\ldots<i_k \leq n}}
    $$
    eine Basis der Clifford-Algebra mit Dimension $\operatorname{dim}_{\mathbb{K}} \mathrm{Cl}(V, \beta)=2^{\operatorname{dim}_{\mathbb{K}} V}$.
  \item \emph{$\Z^2$-Graduierung der Clifford Algebra und kanonischer Automorphismus}: Sei $\imath: V \rightarrow \mathrm{Cl}(V, \beta)$ die Standardeinbettung und $\imath^{\prime}: V \rightarrow \mathrm{Cl}(V, \beta)$ definiert durch
    $$
    \imath^{\prime}(v)=-\imath(v), \quad  v \in V .
    $$
    Da $\imath^{\prime}(v)^2=\imath(v)^2=-\beta(v, v) \cdot 1$ gilt, existiert wegen universeller Eigenschaft der eindeutige \emph{kanonische Algebren-Automorphismus} $\phi: \mathrm{Cl}(V, \beta) \rightarrow \mathrm{Cl}(V, \beta)$, so dass das Diagramm
    \[
    \begin{tikzcd}
    V \arrow[r, "\iota"] \arrow[dr, "\iota'"'] & \Cl(V,\beta) \arrow[d, "\phi"] \\
    & \Cl(V,\beta)
    \end{tikzcd}
    \]
    kommutiert. Durch Identifikation $\imath(V) \subset \mathrm{Cl}(V, \beta)$ mit $V$ haben wir $\left.\phi\right|_V=-\mathrm{id}$. Das induziert eine Dekomposition 
    $$
    \mathrm{Cl}(V, \beta)=\mathrm{Cl}^0(V, \beta) \oplus \mathrm{Cl}^1(V, \beta),
    $$
    mit
    $$
    \begin{aligned}
    & \mathrm{Cl}^0(V, \beta)=+1-\text {Eigenraum von } \phi \\
    & \mathrm{Cl}^1(V, \beta)=-1-\text {Eigenraum von } \phi.
    \end{aligned}
    $$
\end{enumerate}
\end{REM}

\begin{DEF}{}{Clifford-Modul}
Sei $\C l(V,\beta)$ eine komplexifizierte Clifford-Algebra einer reellen Clifford-Algebra. Ein \emph{Clifford-Modul} über $\C l(V,\beta)$ ist ein komplexer Vektorraum $W$ zusammen mit einer Algebrenabbildung
\[
\rho:\C l(V,\beta)\longrightarrow \End_{\C}(W).
\]
oder äquivalent einer \emph{Wirkung der Clifford-Algebra auf $W$}:
\[
\C l(V,\beta)\times W\longrightarrow W,\qquad
(z,\psi)\longmapsto z\cdot \psi:=\rho(z)(\psi)
\]
Eine Clifford-Moduln-Struktur induziert eine \emph{Clifford-Multiplikation} auf $W$ über
$\C l(V,\beta)$:
\[
c:=\rho|_V:V\to \End_{\C}(W)
\]
und man schreibt
\[
c(v)\psi=v\cdot\psi \qquad v\in V,\ \psi\in W.
\]
Die Clifford-Relation lautet dann
\[
c(v)c(w)+c(w)c(v)=-2\,g(v,w)\,\id_{W}.
\]
Wegen der Clifford-Relation bestimmt $c$ mittels der universellen Eigenschaft von $\C l(V, \beta)$ eindeutig die Darstellung $\rho$.

\end{DEF}


\begin{DEF}{}{Pin- und Spin-Gruppe}
\begin{enumerate}
\item Die \emph{Pin-Gruppe} ist definiert als die von den Einheitsvektoren erzeugte Untergruppe
\[
\Pin(n):=\langle S^{n-1}\rangle = \left\{v_1 \cdot \ldots \cdot v_m \in \mathrm{Cl}_n \mid v_j \in S^{n-1} \subset \mathbb{R}^n, m \in \mathbb{N}_0\right\} \subset \Cl_n^\times.
\]

\item Die \emph{Spin-Gruppe} ist der gerade Teil von der Pin Gruppe:
$$
\begin{aligned}
\operatorname{Spin}(n) & :=\operatorname{Pin}(n) \cap \mathrm{Cl}_n^0 \\
& =\left\{v_1 \cdot \ldots \cdot v_m \in \mathrm{Cl}_n \mid v_j \in S^{n-1}, m \in 2 \mathbb{N}_0\right\} .
\end{aligned}
$$
\end{enumerate}
\end{DEF}


\begin{REM}{}{Eigenschaften der Spin Gruppe}
  \begin{enumerate}
    \item \emph{Lie-Untergruppe}: $\Spin(n)$ ist eine abgeschlossene Lie-Untergruppe von $\Cl_n^\times$ und $\Pin(n)$.
    \item \emph{Zweifache Überlagerung}: Man definiert durch  
    \[
    \Lambda:\Pin(n)\to \mathrm{O}(n);\ u\mapsto [x\mapsto \Lambda(u)(x):=\alpha(u)\,x\,u^{-1}]
    \]
    ein surjektiver Gruppenhomomorphismus mit $\ker\Lambda=\{\pm 1\}$, wobei $\alpha:\Cl_n\to\Cl_n$ die Gradinvolution ist, d.h.
    \[
    \alpha(v)=-v \quad \text{für } v\in \R^n,
    \qquad
    \alpha(x)=(-1)^k x \text{ für } x\in \Cl_n^{(k)}.
    \]
    Einschränkung von $\Lambda$ auf die Spin Gruppe $\Spin(n)\subset\Pin(n)$ liefert einen surjektiven Liegruppenhomomorphismus
    \[
    \lambda:=\Lambda|_{\Spin(n)}:\Spin(n)\to \SO(n),
    \qquad
    \lambda(u)(x)=u x u^{-1},
    \]
    auch mit $\ker\lambda=\{\pm 1\}$, die kanonische zweifache Überlagerung von $\SO(n)$ durch $\Spin(n)$.
  \end{enumerate}
\end{REM}




\end{document}